\fancypagestyle{plain}{
	\fancyhf{}
	\fancyfoot[R]{\thepage}
	\renewcommand{\headrulewidth}{0pt}
	\renewcommand{\footrulewidth}{0pt}
}

\chapter{Introduction}
\label{chapter:introduction}

\section{Problem Statement}

Hiring used to run on paper. Job seekers printed resumes and delivered them in person or by post, and employers advertised openings in newspapers and reviewed applications by hand. Online recruitment platforms have changed this: employers now publish vacancies and candidates apply electronically, so distance and response time no longer decide how fast a role is filled.

The new approach brings its own friction. Each recruitment site has its own search interface, its own resume format, and its own application flow. A candidate who wants to reach more than one pool of jobs must recreate essentially the same profile and resume on every platform, then visit each one separately to check whether an application has progressed. The more platforms a job seeker uses, the more duplicated effort the search costs.

Employers face a similar mess in reverse. Publishing a posting is only the first step; they still have to collect applications, read resumes, record decisions, and keep those records straight. Without a single place that ties postings, applications, and decisions together, the employer ends up stitching information across tools — recreating the inefficiency the platform was supposed to remove.

These problems motivate a centralized system: one platform where job seekers search jobs, maintain a profile, upload a resume, apply, and track their applications, and where employers manage company information, publish postings, and review applicants through one interface. Because such a system stores personal data and lets users modify shared records, access must be controlled. Registration and login protect each user's data, and role-based permissions decide what each user may do: job seekers manage their profiles and applications, employers manage postings and candidates, and a back-office role keeps platform content in check. This thesis designs and implements that system, focusing on two central flows — applying for jobs on the job seeker's side, and managing postings and their applicants on the employer's side.